\documentclass[11pt,a4paper]{article}
\usepackage[utf8]{inputenc}
\usepackage[T1]{fontenc}
\usepackage{amsmath,amssymb}
\usepackage[margin=1in]{geometry}
\usepackage{natbib}
\usepackage{url}
\usepackage{hyperref}

\newcommand{\provisional}{\textit{[Provisional: notation and functional forms subject to change.]}}

\title{Rates, Runoff and Sales: Optimal Monetary and Balance-Sheet Policy\\ in a Nonlinear HANK Model with a Preferred-Habitat Gilt Market\thanks{Working outline. All content is preliminary.}}
\author{Cyrille Lenoel}
\date{Draft outline, \today}

\begin{document}
\maketitle

\begin{abstract}
\noindent [To be written.] This paper studies the joint conduct of interest-rate and balance-sheet policy in a nonlinear New Keynesian model with heterogeneous households and a preferred-habitat gilt market, calibrated to the United Kingdom. Gilt purchases, passive run-off and active sales are modelled as distinct instruments. The model is solved globally with deep neural networks, verified against the linearised solution, and disciplined by high-frequency evidence from UK gilt markets.
\end{abstract}

\section{Introduction}
\begin{itemize}
  \item \textbf{Motivation.} Since 2009 the Bank of England has used both Bank Rate and the Asset Purchase Facility (APF). Since 2022 it has reduced the APF through passive run-off \emph{and} outright gilt sales, which sets it apart from the Federal Reserve. The September--November 2022 LDI episode showed that the gilt market can be fragile when pension and LDI demand for long-dated gilts is concentrated.
  \item \textbf{Question.} How should the short rate and the three balance-sheet instruments (QE, passive QT and active QT) be combined when term premia respond nonlinearly to the size and composition of central-bank holdings? Does the policy split in \citet{KamdarRay2026} still hold in this setting? In that split, the short rate stabilises output and inflation and the balance sheet stabilises term premia.
  \item \textbf{Approach.} We build a HANK model with preferred-habitat arbitrageurs à la \citet{VayanosVila2021}. Habitat demand is calibrated to UK pension and LDI holdings, separately for conventional and index-linked gilts \citep{ChenBeetsmaVanWijnbergen}. The model is solved globally with deep equilibrium nets \citep{AzinovicGaegaufScheidegger2022, MaliarMaliarWinant2021}.
  \item \textbf{Contributions.} (i) A model in which QE, passive QT and active QT are separate instruments with potentially asymmetric effects. (ii) A global deep-learning solution that is verified systematically against the linearised solution wherever both apply. (iii) Empirical discipline from UK high-frequency gilt local projections, including a post-2022 regime interaction and the LDI episode.
\end{itemize}

\section{Related Literature}
\noindent\textit{Note: this section will be filled from EconResearcher's literature review in \texttt{dl-jmp/literature/}.}
\begin{itemize}
  \item \textbf{Preferred habitat and QE/QT.} Term-structure models with segmented demand \citep{VayanosVila2021}; optimal monetary and balance-sheet policy with preferred habitat \citep{KamdarRay2026}; UK pension and LDI demand across conventional and index-linked gilts \citep{ChenBeetsmaVanWijnbergen}; fiscal backing and central-bank net worth \citep{GrodeckaMessiKliemMuller}.
  \item \textbf{Deep-learning solution methods.} \citet{MaliarMaliarWinant2021}; \citet{AzinovicGaegaufScheidegger2022}; \citet{FernandezVillaverdeHurtadoNuno2023}; \citet{EbrahimiKahou2021, EbrahimiKahou2024}; \citet{Kase2024}; \citet{Nuno2024}; \citet{Duarte2024}. See the overview in \citet{FVDeckHarvard}.
  \item \textbf{High-frequency identification of monetary policy.} Separating policy from information shocks with sign restrictions \citep{JarocinskiKaradi2020}; UK intraday surprise databases \citep{BraunMirandaAgrippinoSaha}.
\end{itemize}

\section{Model}
\provisional

\subsection{Households}
A continuum of households faces idiosyncratic income risk $z_{it}$ and a borrowing limit, and saves in a liquid short asset:
\begin{align}
  V(a,z;\mathbf{S}) &= \max_{c,a'}\; u(c) + \beta\,\mathbb{E}\left[V(a',z';\mathbf{S}')\right] \\
  \text{s.t.}\quad c + a' &= (1+r)\,a + (1-\tau)\,w\,z\,n + T, \qquad a' \ge \underline{a}.
\end{align}
The aggregate state $\mathbf{S}$ includes the wealth distribution $\mu$, which is approximated by a finite set of moments or by a permutation-invariant representation \citep{EbrahimiKahou2021}.

\subsection{Preferred-Habitat Arbitrageurs}
Risk-averse arbitrageurs with risk aversion $a$ trade gilts of maturity $\tau$ against habitat investors, whose demand is
\begin{equation}
  Z_t(\tau) = -\alpha(\tau)\left[\log P_t(\tau) - \beta(\tau)\right] + \theta_t(\tau),
\end{equation}
where $\theta_t(\tau)$ captures pension and LDI demand shocks. Market clearing requires
\begin{equation}
  X_t(\tau) + Z_t(\tau) + B^{CB}_t(\tau) = B^{G}_t(\tau),
\end{equation}
where $B^{CB}_t(\tau)$ is the stock of central-bank holdings. Term premia therefore depend on the size and maturity composition of $B^{CB}_t$ and on how much risk-bearing capacity arbitrageurs have.

\subsection{Firms}
Monopolistically competitive firms face Rotemberg price-adjustment costs, which give a nonlinear Phillips curve:
\begin{equation}
  \left(\pi_t - \bar\pi\right)\pi_t = \frac{\epsilon}{\psi}\left(mc_t - \frac{\epsilon-1}{\epsilon}\right) + \mathbb{E}_t\left[\Lambda_{t,t+1}\frac{Y_{t+1}}{Y_t}\left(\pi_{t+1}-\bar\pi\right)\pi_{t+1}\right].
\end{equation}

\subsection{Policy Block: Three Balance-Sheet Instruments}
The short rate follows a rule subject to an effective lower bound:
\begin{equation}
  i_t = \max\left\{\underline{i},\; \bar{i} + \phi_\pi(\pi_t-\bar\pi) + \phi_y \hat{y}_t + \varepsilon^i_t\right\}.
\end{equation}
Central-bank holdings evolve as
\begin{equation}
  B^{CB}_{t+1}(\tau) = \underbrace{(1-\delta(\tau))\,B^{CB}_t(\tau)}_{\text{passive run-off}} + \underbrace{Q^{+}_t(\tau)}_{\text{QE}} - \underbrace{Q^{-}_t(\tau)}_{\text{active sales}}, \qquad Q^{\pm}_t(\tau)\ge 0.
\end{equation}
Active sales may carry a market-impact cost $\kappa\,(Q^{-}_t)^2$, which creates an asymmetry with QE. A consolidated budget links the APF to the Treasury through an indemnity regime (Section~\ref{sec:ext}).

\section{Solution Method}
\begin{itemize}
  \item \textbf{Network.} A feed-forward network $\mathcal{N}_\rho(\mathbf{S})$ maps the aggregate state, including the distribution summary and the policy state $B^{CB}$, to policy functions and prices. Kinks such as the lower bound and $Q^{\pm}\ge0$ are handled with Fischer--Burmeister complementarity terms \citep{MaliarMaliarWinant2021, AzinovicGaegaufScheidegger2022}.
  \item \textbf{Loss.} The loss is built from equilibrium residuals (Euler equations, the Phillips curve and gilt market clearing). Conditional expectations are formed from \emph{two independent shock draws}, so that the product of the residuals is an unbiased estimate of the squared expected residual (the all-in-one approach of \citealp{MaliarMaliarWinant2021}).
  \item \textbf{Training.} States are drawn from the model's ergodic set by simulation. The network is first trained with Adam and then refined with L-BFGS to reach a low residual floor.
  \item \textbf{Diagnostics.} We report residual distributions on held-out ergodic draws, check robustness to architecture and seeds, and track the accuracy of aggregation.
\end{itemize}

\section{Verification against the Linear Solution}
\begin{itemize}
  \item \textbf{Rationale.} A deep-learning solution comes with no general guarantee of correctness. We therefore compare it with the linearised solution wherever both apply.
  \item \textbf{Nested cases.} Small shocks around the steady state, the representative-agent limit and a parameterisation without a lower bound should all reproduce the first-order impulse responses and the linear policy rules.
  \item \textbf{Metrics.} The maximum and mean deviations of impulse responses and policy functions, reported over a grid of shock sizes, show where nonlinearity starts to matter.
  \item \textbf{Divergence analysis.} Gaps between the global and linear solutions away from the steady state are attributed to specific nonlinearities: the lower bound, arbitrageur risk-bearing capacity and the asymmetry between sales and purchases.
\end{itemize}

\section{Empirical Targets}
\begin{itemize}
  \item \textbf{Sample.} 175 MPC announcement days from 2009 to September 2026, using the team's UK high-frequency gilt local projections.
  \item \textbf{Shock purification.} A Jaroci\'nski--Karadi rate--equity sign filter \citep{JarocinskiKaradi2020} keeps 91 policy days on the FTSE 100, and 71 days are robust under both the FTSE 100 and FTSE 250.
  \item \textbf{Specification.} Local projections of gilt yields across maturities include a post-2022 interaction and a dummy for the September--November 2022 LDI episode. Surprises are benchmarked against the Braun--Miranda-Agrippino--Saha intraday surprises \citep{BraunMirandaAgrippinoSaha}.
  \item \textbf{Use in the model.} The model is calibrated or estimated by matching its impulse responses to these, with neural-network estimation \citep{Kase2024} as one option.
\end{itemize}

\section{Planned Results}
\begin{itemize}
  \item Whether the split in \citet{KamdarRay2026} (the short rate for output and inflation, the balance sheet for term premia) survives household heterogeneity and global nonlinearity.
  \item How the effects of QE, passive QT and active QT on term premia and on activity differ, and in particular whether these effects depend on the state.
  \item How fragile the market is when LDI demand is concentrated, and what that implies for the pace and maturity of active sales.
  \item The distributional incidence of balance-sheet policy across the wealth distribution.
\end{itemize}

\section{Extensions}\label{sec:ext}
\begin{itemize}
  \item \textbf{Fiscal backing and central-bank net worth.} APF losses are modelled under full, delayed or no Treasury indemnity \citep{GrodeckaMessiKliemMuller}.
  \item \textbf{Yield curve control.} We compare quantity-based QT with price-based targets and ask whether the long rate is left free to move \citep{FukuiYagasaki}.
  \item \textbf{Fiscal delay.} Delays in the fiscal response, following \citet{BayerGongKuester}, interact with balance-sheet policy.
  \item \textbf{Fallbacks.} (a) A global solution of a sectoral input--output NK model with a binding lower bound. (b) Neural-network estimation of a nonlinear UK model following \citet{Kase2024}.
\end{itemize}

\section{Conclusion}
\begin{itemize}
  \item This section will summarise the answer to the policy-assignment question for the UK.
  \item It will draw out the lessons for the design of QT, in particular the choice between run-off and sales.
  \item It will set out the methodological lesson: global deep-learning solutions are credible when they are verified against tractable benchmarks.
\end{itemize}

\bibliographystyle{plainnat}
\bibliography{refs}
\end{document}
